% ============================================================
%  Aegis Entropy — Academic Project Report
%  AI-Based Intelligent System for Network Security
%  Monitoring & Automated Response
%  Module: Artificial Intelligence | ENSA Kénitra
%  Academic Year 2025–2026
% ============================================================
\documentclass[12pt,a4paper]{report}

\usepackage[utf8]{inputenc}
\usepackage[T1]{fontenc}
\usepackage[english]{babel}
\usepackage{geometry}
\usepackage{booktabs}
\usepackage{array}
\usepackage{tabularx}
\usepackage{longtable}
\usepackage{multirow}
\usepackage{enumitem}
\usepackage{fancyhdr}
\usepackage{titlesec}
\usepackage{hyperref}
\usepackage{parskip}
%\usepackage{microtype}
\usepackage{amssymb}
\usepackage{setspace}

\geometry{margin=2.5cm, top=3cm, bottom=3cm}

% --- Column types ---
\newcolumntype{L}[1]{>{\raggedright\arraybackslash}p{#1}}
\newcolumntype{C}[1]{>{\centering\arraybackslash}p{#1}}
\newcolumntype{B}[1]{>{\bfseries\raggedright\arraybackslash}p{#1}}

% --- Section formatting ---
\titleformat{\chapter}[display]
  {\normalfont\Large\bfseries}{\chaptername\ \thechapter}{12pt}{\LARGE}
\titlespacing{\chapter}{0pt}{0pt}{20pt}

\titleformat{\section}
  {\normalfont\large\bfseries}{\thesection}{1em}{}
\titleformat{\subsection}
  {\normalfont\normalsize\bfseries}{\thesubsection}{1em}{}

% --- Header/Footer ---
\pagestyle{fancy}
\fancyhf{}
\fancyhead[L]{\small Aegis Entropy --- Academic Project Report \quad Module: Artificial Intelligence}
\fancyhead[R]{\small ENSA K\'{e}nitra --- Networks \& Telecommunication Systems}
\fancyfoot[C]{\small Page \thepage}
\renewcommand{\headrulewidth}{0.4pt}

% --- Hyperref ---
\hypersetup{colorlinks=false, hidelinks}

% --- Spacing ---
\setlength{\parskip}{6pt}
\setlength{\parindent}{0pt}
\setlength{\tabcolsep}{6pt}

% ============================================================
\begin{document}

% ============================================================
%  TITLE PAGE
% ============================================================
\begin{titlepage}
  \centering
  \vspace*{2cm}
  {\Huge\bfseries Aegis Entropy}\\[0.8em]
  {\Large AI-Based Intelligent System for Network Security\\
  Monitoring \& Automated Response}\\[1.2em]
  \rule{\textwidth}{0.8pt}\\[0.4em]
  {\large\textsc{Port Scanning Detection \(\bullet\) MTD \(\bullet\) Honeypot}}\\[0.4em]
  \rule{\textwidth}{0.4pt}
  \vspace{1.5cm}

  \begin{tabular}{@{}ll@{}}
    \textbf{Module:}        & Artificial Intelligence \\[4pt]
    \textbf{Status:}        & Academic Project \\[4pt]
    \textbf{Severity:}      & MEDIUM-HIGH \\[4pt]
    \textbf{Academic Year:} & 2025 -- 2026 \\
  \end{tabular}

  \vspace{1.8cm}
  {\large\bfseries Project Team}\\[0.8em]
  \begin{tabular}{@{}c@{}}
    El Kartouti Anas \\[3pt]
    Hammoubel El Yazid \\[3pt]
    Lekhbioui Adil \\[3pt]
    Moulay Anas \\[3pt]
    Nafia Khadija \\
  \end{tabular}

  \vfill
  {\normalsize ENSA K\'{e}nitra --- Networks \& Telecommunication Systems}
\end{titlepage}

% ============================================================
%  TABLE OF CONTENTS
% ============================================================
\tableofcontents
\clearpage

% ============================================================
%  PART I — NEEDS EXPRESSION
% ============================================================
\part{Needs Expression}

\textit{Document origin: This part is derived from the initial Needs Expression document
(Aegis-Morph phase) which establishes the academic, methodological, and contextual
foundations of the Aegis Entropy project.}

% ------------------------------------------------------------
\chapter{Project Contextualisation}

\section{General Context}

The digital transformation of organisations across all sectors --- public administration,
healthcare, finance, industry, and education --- has led to an unprecedented expansion of
networked infrastructures. Whether deployed as Local Area Networks (LAN), Wide Area
Networks (WAN), Metropolitan Area Networks (MAN), wireless networks, cloud environments,
or hybrid architectures, these infrastructures now serve as the backbone of critical operations,
enabling resource sharing, real-time communication, and distributed service delivery.

This ubiquity, however, comes at a significant cost: the larger and more complex a network,
the broader its attack surface. Network security has consequently become one of the most
pressing challenges in modern information systems management. Malicious actors continuously
evolve their strategies, exploiting vulnerabilities at every layer of the network stack --- from the
physical and data-link layers up through the application layer.

In response, the field of Network Intrusion Detection and Prevention Systems (IDS/IPS) has
emerged as a critical discipline. Traditional approaches rely on predefined signatures and static
rule sets. While effective against known threats, these systems fail to anticipate novel,
adaptive, or previously unseen attack patterns. The limitations of such approaches have
created a clear and growing demand for intelligent, self-learning security systems capable of
analysing network behaviour dynamically and responding autonomously.

% ------------------------------------------------------------
\chapter{The Role of Artificial Intelligence in Network Security}

Artificial Intelligence, and more specifically machine learning, offers a fundamentally different
approach to intrusion detection. Rather than relying on manually curated rules, AI-based
systems learn statistical models of normal network behaviour from labelled traffic datasets.
Any significant deviation from the learned baseline is flagged as a potential threat --- enabling
the detection of both known and previously uncharacterised attack vectors.

This paradigm shift brings several decisive advantages over classical IDS/IPS solutions:
adaptability to evolving threats, the ability to process and classify high-volume traffic in real
time, reduced dependency on manual signature updates, and the capacity to support
autonomous decision-making --- including automated threat containment without human
intervention.

\begin{table}[h!]
\centering
\caption{Classical IDS/IPS vs.\ AI-Based IDS/IPS}
\begin{tabular}{@{}L{7cm} L{7cm}@{}}
  \toprule
  \textbf{Classical IDS/IPS} & \textbf{AI-Based IDS/IPS} \\
  \midrule
  Relies on predefined signatures          & Learns behavioural patterns from data \\
  \addlinespace
  Fails against unknown/novel attacks      & Detects anomalies beyond known signatures \\
  \addlinespace
  Requires constant manual updates         & Self-adapts through model retraining \\
  \addlinespace
  No autonomous response capability        & Supports automated threat neutralisation \\
  \addlinespace
  High false positive rate at scale        & Optimised precision/recall via ML tuning \\
  \addlinespace
  Static and network-topology-specific     & Generalisable across diverse network types \\
  \bottomrule
\end{tabular}
\label{tab:ids-comparison}
\end{table}

% ------------------------------------------------------------
\chapter{Project Framework \& Methodology}

This project is carried out within the Artificial Intelligence academic module by a team of
five students. Its objective is to design, implement, test, and validate an AI-powered IDS/IPS
system, deployed over a simulated network topology spanning multiple nodes and traffic types.
The system must be capable of monitoring network traffic in real time, detecting suspicious
behavioural patterns, profiling threat actors, triggering alerts, and enforcing automated
countermeasures.

The project is structured around the \textbf{PDCA (Plan--Do--Check--Act)} continuous
improvement methodology, which ensures that each development phase is rigorously planned,
executed, evaluated, and refined before progression to the next stage.

\begin{table}[h!]
\centering
\caption{PDCA Methodology Phases}
\begin{tabular}{@{}B{2cm} L{12cm}@{}}
  \toprule
  \textbf{Phase} & \textbf{Description} \\
  \midrule
  Plan  & Define project objectives, identify data sources, design the network topology,
          select AI models and evaluation criteria, and structure the development roadmap. \\
  \addlinespace
  Do    & Implement the full system pipeline: data preprocessing, model training, real-time
          traffic capture, threat detection engine, dashboard, and automated response module. \\
  \addlinespace
  Check & Evaluate the system against simulated threat scenarios. Measure performance
          against predefined metrics (accuracy, latency, false positive rate). Document
          observations and gaps. \\
  \addlinespace
  Act   & Validate findings through cross-validation and end-to-end testing. Identify
          limitations, propose improvements, and consolidate results into the final academic
          deliverable. \\
  \bottomrule
\end{tabular}
\label{tab:pdca}
\end{table}

% ------------------------------------------------------------
\chapter{Design Methodology: CPS}

Beyond the overarching PDCA framework, the project adopts the
\textbf{CPS (Context--Problematic--Solution)} scientific design methodology to ensure that
every technical decision is grounded in observed, field-validated needs rather than
assumptions.

\begin{table}[h!]
\centering
\caption{CPS Methodology Phases}
\begin{tabular}{@{}B{2.5cm} L{11.5cm}@{}}
  \toprule
  \textbf{Phase} & \textbf{Description} \\
  \midrule
  Context     & Establishes the environment in which the project operates: the nature of
                modern network infrastructures, the evolving threat landscape, and the
                limitations of existing security tools. This Needs Expression document
                constitutes the Context phase. \\
  \addlinespace
  Problematic & Derived from structured interviews conducted with network and infrastructure
                professionals. An interview guide will be developed, administered, and
                transcribed into formal interview reports. The technical problematic will be
                formalised exclusively from the insights gathered during this phase. \\
  \addlinespace
  Solution    & A targeted technical solution is proposed based on the findings of the
                Problematic phase. It specifies the chosen approach, tools, datasets, model
                architecture, evaluation metrics, and the integration of advanced defence
                mechanisms into the IDS/IPS system. \\
  \bottomrule
\end{tabular}
\label{tab:cps}
\end{table}

% ------------------------------------------------------------
\chapter{Network Scope}

The system is designed to be applicable across a broad spectrum of network architectures and
deployment contexts. While the academic prototype will be validated in a controlled simulation
environment, the principles, models, and mechanisms developed are intended to generalise to
real-world infrastructure of any type and scale.

\begin{table}[h!]
\centering
\caption{Applicable Network Types}
\begin{tabular}{@{}L{3cm} L{5.5cm} L{5.5cm}@{}}
  \toprule
  \textbf{Network Type} & \textbf{Description} & \textbf{Relevance} \\
  \midrule
  LAN
    & Local Area Network --- controlled, high-speed, limited geographic scope
    & Primary simulation environment for prototype validation \\
  \addlinespace
  WAN
    & Wide Area Network --- geographically distributed, multi-site connectivity
    & Relevant for enterprise and inter-organisational threat detection \\
  \addlinespace
  MAN
    & Metropolitan Area Network --- city-scale infrastructure
    & Applicable to smart city, utility, and campus-wide deployments \\
  \addlinespace
  Wireless / Wi-Fi
    & IEEE 802.11 networks --- mobile, access-point-based
    & Relevant to BYOD environments and open-access networks \\
  \addlinespace
  Cloud / Virtual
    & Software-defined and cloud-hosted network infrastructure
    & Growing target surface for network-level intrusions \\
  \addlinespace
  Hybrid
    & Combination of on-premise and cloud, wired and wireless
    & Reflects most real-world organisational deployments \\
  \bottomrule
\end{tabular}
\label{tab:network-scope}
\end{table}

% ------------------------------------------------------------
\chapter{Advanced Security Concepts Integrated in the Project}

\section{Port Polymorphism via Aegis-Morph}

Aegis-Morph is a dynamic port reassignment mechanism that continuously rotates the
listening ports of protected network services according to a defined schedule or trigger
condition. By making service port assignments unpredictable and transient, this technique
disrupts an attacker's ability to map the network and identify exploitable entry points through
conventional probing strategies. The AI detection engine is designed with full awareness of the
current port assignment state, ensuring that legitimate traffic to dynamically assigned ports is
not misclassified as suspicious.

\section{Honeypot Integration}

A honeypot is a deliberately exposed decoy system deployed within the network
infrastructure. It mimics the appearance of a legitimate host or service while containing no
actual operational data or functionality. Any interaction with the honeypot is, by definition,
unauthorised and constitutes a high-confidence threat signal. Honeypot-captured events enrich
the AI engine's threat intelligence, provide ground-truth labelled data for attacker behavioural
analysis, and enable the generation of comprehensive attacker profiles including intent,
methodology, and originating context.

% ------------------------------------------------------------
\chapter{Expected System Capabilities}

\begin{table}[h!]
\centering
\caption{Expected System Capabilities}
\begin{tabular}{@{}L{5.5cm} L{9cm}@{}}
  \toprule
  \textbf{Capability} & \textbf{Description} \\
  \midrule
  Real-Time Traffic Monitoring
    & Continuous capture and analysis of network flows with live dashboard indicators. \\
  \addlinespace
  AI-Driven Threat Detection
    & Machine learning model trained on labelled network traffic datasets to identify
      abnormal behaviours. \\
  \addlinespace
  Attacker Profiling
    & Automatic generation of a threat actor profile (network identity, behaviour pattern,
      threat level). \\
  \addlinespace
  Graduated Alert System
    & Multi-level alerting mechanism triggered proportionally to confidence and severity
      of detected threats. \\
  \addlinespace
  Automated Blocking
    & Autonomous enforcement of firewall countermeasures upon confirmed threat
      detection, without human intervention. \\
  \addlinespace
  Port Polymorphism (Aegis-Morph)
    & Dynamic rotation of service port assignments to increase the cost and complexity
      of network reconnaissance. \\
  \addlinespace
  Honeypot-Based Deception
    & Decoy node that captures unauthorised interaction attempts and feeds behavioural
      data to the AI engine. \\
  \addlinespace
  Cross-Network Applicability
    & Architecture designed for generalisation across LAN, WAN, wireless, cloud, and
      hybrid environments. \\
  \bottomrule
\end{tabular}
\label{tab:capabilities}
\end{table}

% ============================================================
%  PART II — DETAILED NEEDS EXPRESSION
% ============================================================
\part{Detailed Needs Expression --- Port Scanning, MTD \& Honeypot}

\textit{Document origin: This part corresponds to the formal, finalised Needs Expression
document (PortScan\_MTD\_Honeypot\_Final), which constitutes the complete technical
decision-making specification for the Aegis Entropy system.}

% ------------------------------------------------------------
\chapter{Executive Summary}

This document presents the formal decision-making needs expression for the design,
implementation, testing, and validation of an Artificial Intelligence model dedicated to the
\textbf{real-time detection and automated mitigation} of Port Scanning and Network
Reconnaissance activity within networked environments.

Reconnaissance constitutes the first phase of virtually every cyberattack kill chain: an
adversary probes the network to map active hosts, enumerate open ports, and fingerprint
running services before launching a targeted intrusion. Early detection and containment at
this phase can prevent the full materialisation of a subsequent, more destructive attack.

Beyond detection, this project integrates two advanced and complementary proactive defence
mechanisms: \textbf{Moving Target Defence (MTD)}, which continuously mutates the
network's observable attack surface to disrupt and mislead scanning activity, and
\textbf{Honeypots}, which deploy decoy assets to attract, trap, and profile adversarial
reconnaissance behaviour. Together, these three pillars --- detection, deception, and surface
mutation --- form a layered, intelligent security posture.

% ------------------------------------------------------------
\chapter{Problem Statement}

Port scanning, in its various forms --- SYN scan, ACK scan, UDP scan, XMAS scan, FIN scan
--- generates traffic patterns that differ fundamentally from legitimate host-to-host
communication. An adversary conducting a scan sends probes to hundreds or thousands of
ports in succession, producing an anomalous distribution of half-open connections, RST
responses, and unanswered packets.

Despite being a well-understood technique, detecting \textbf{stealthy or slow-rate scanning}
--- where the adversary deliberately spaces probes to evade threshold-based detection ---
remains a significant challenge. Such low-and-slow patterns require behavioural analysis over
extended time windows rather than instantaneous packet-level inspection.

A further critical dimension is the \textbf{static, predictable nature of conventional network
configurations}. Fixed IP addresses, stable port assignments, and consistent service layouts
give adversaries a reliable map to exploit. Once a scan is complete, the attacker possesses a
stable, actionable picture of the infrastructure --- a problem that no detection system alone
can fully address.

% ------------------------------------------------------------
\chapter{Project Objectives}

\begin{table}[h!]
\centering
\caption{Project Objectives}
\begin{tabular}{@{}L{4.5cm} L{10cm}@{}}
  \toprule
  \textbf{Objective} & \textbf{Description} \\
  \midrule
  Primary Objective
    & Train and deploy an AI model capable of detecting port scanning and network
      reconnaissance activity in real time across networked environments. \\
  \addlinespace
  Scan Type Coverage
    & Detect and classify SYN Scan (stealth), TCP Connect Scan, UDP Scan, ACK Scan,
      XMAS Scan, and slow/distributed scanning variants. \\
  \addlinespace
  Moving Target Defence
    & Integrate an MTD mechanism that continuously rotates network configuration
      parameters --- including port assignments and IP addresses --- to invalidate
      adversarial reconnaissance maps. \\
  \addlinespace
  Honeypot Deception
    & Deploy honeypot nodes that mimic legitimate services to attract, trap, and observe
      adversarial scanning activity, enriching threat intelligence. \\
  \addlinespace
  Attacker Profiling
    & Automatically build a structured profile for each detected threat actor: network
      identity, scan methodology, targeted scope, behavioural timeline. \\
  \addlinespace
  Alerting
    & Trigger real-time, graduated alerts upon detection of scanning activity, classified by
      scan type, confidence score, and severity level. \\
  \addlinespace
  Automated Response
    & Enforce automated protective measures --- including network-level blocking --- upon
      confirmed threat detection, without requiring human intervention. \\
  \bottomrule
\end{tabular}
\label{tab:objectives}
\end{table}

% ------------------------------------------------------------
\chapter{Functional Requirements}

\section{System Inputs}

\begin{itemize}[nosep, leftmargin=1.5em]
  \item Raw TCP/UDP packets captured on the network interface, filtered for partial or
        incomplete handshake patterns indicative of probing behaviour.
  \item Labelled training dataset containing port scanning and benign traffic samples
        (CICIDS2017 PortScan subset, UNSW-NB15 Reconnaissance category).
  \item Sliding time-window aggregated flow features computed per source IP over
        configurable observation windows.
  \item Network metadata: source/destination IP, MAC addresses, ports, protocol, TCP
        flag combinations, TTL values, and timestamps.
  \item Honeypot interaction logs: connection attempts, payloads, source identifiers, and
        session metadata from decoy nodes.
  \item MTD state feed: current port assignment table and IP rotation schedule, enabling
        the AI engine to correctly interpret traffic against the dynamic topology.
\end{itemize}

\section{System Outputs}

\begin{itemize}[nosep, leftmargin=1.5em]
  \item Real-time dashboard: per-source port probe counts, scan heat map, active threat
        list, honeypot interaction events, MTD rotation log.
  \item Per-time-window classification: \textsc{Benign} or \textsc{Port\_Scan}, with
        inferred scan type and associated confidence score.
  \item Attacker profile card: source IP and MAC, scan type, ports and hosts targeted,
        first/last seen timestamps, honeypot interaction flag, threat level.
  \item Graduated alert notification with scan type label, severity score, and contextual
        metadata.
  \item Automated network-level block rule applied to the confirmed attacker's source
        address, with the event logged to the incident database.
  \item MTD rotation trigger: upon detection, the system initiates a port/IP rotation cycle
        to further invalidate any reconnaissance data collected by the attacker.
\end{itemize}

% ------------------------------------------------------------
\chapter{Dataset Specification}

\begin{table}[h!]
\centering
\caption{Training Dataset Specification}
\begin{tabular}{@{}L{3.5cm} L{4.5cm} C{3cm} L{3.5cm}@{}}
  \toprule
  \textbf{Dataset} & \textbf{Relevant Subset} & \textbf{Approx.\ Samples} & \textbf{Source} \\
  \midrule
  CICIDS2017   & PortScan category             & $\sim$158,000 flows  & UNB / CIC \\
  UNSW-NB15    & Reconnaissance category       & $\sim$13,000 records & UNSW Australia \\
  CAIDA Telescope & Backscatter / SYN flood probes & Variable captures & CAIDA.org \\
  \bottomrule
\end{tabular}
\label{tab:datasets}
\end{table}

% ------------------------------------------------------------
\chapter{Key Network Features for the ML Model}

\begin{table}[h!]
\centering
\caption{Key Network Features for the ML Model}
\begin{tabular}{@{}L{4cm} L{5.5cm} L{4.5cm}@{}}
  \toprule
  \textbf{Feature Name} & \textbf{Description} & \textbf{Relevance} \\
  \midrule
  Distinct Dst Ports/IP
    & \# unique destination ports contacted by source
    & Critical --- core scanning signature \\
  \addlinespace
  SYN Flag Count
    & SYN packets sent without completing handshake
    & Critical --- stealth SYN scan indicator \\
  \addlinespace
  RST Flag Count
    & RST responses from closed ports
    & Very High --- port probe rejection \\
  \addlinespace
  Flow Duration
    & Duration of each individual probe connection
    & High --- very short in fast scans \\
  \addlinespace
  Total Fwd Packets
    & Packets per flow from the scanner
    & High --- usually 1 per probe \\
  \addlinespace
  ACK Flag Count
    & ACK packets sent without a prior SYN
    & Medium --- ACK scan detection \\
  \addlinespace
  IAT Mean
    & Mean inter-arrival time between probes
    & High --- slow scans have large IAT \\
  \addlinespace
  Unique Dst IPs/Src
    & \# distinct destination hosts contacted
    & Medium --- host discovery phase \\
  \addlinespace
  TTL Value
    & Time-to-Live field in the IP header
    & Low --- OS fingerprinting attempts \\
  \addlinespace
  Bwd Packet Length
    & Response packet size
    & Medium --- RST vs SYN-ACK ratio \\
  \addlinespace
  Honeypot Interaction Flag
    & Binary: did source IP contact a decoy node?
    & Critical --- high-confidence threat signal \\
  \addlinespace
  MTD Port Delta
    & Difference between probed port and current assigned port
    & High --- reveals stale reconnaissance data \\
  \bottomrule
\end{tabular}
\label{tab:features}
\end{table}

% ------------------------------------------------------------
\chapter{Machine Learning Model Strategy}

\begin{table}[h!]
\centering
\caption{Machine Learning Model Strategy}
\begin{tabular}{@{}L{4cm} L{10.5cm}@{}}
  \toprule
  \textbf{Component} & \textbf{Details} \\
  \midrule
  Primary Model
    & Random Forest Classifier --- handles high-dimensional scan features with natural
      feature importance ranking and robustness to class imbalance. \\
  \addlinespace
  Secondary Model
    & XGBoost --- compared for accuracy on time-windowed aggregated features and as
      a high-performance ensemble benchmark. \\
  \addlinespace
  Slow Scan Detection
    & Isolation Forest (unsupervised) applied on sliding 60-second window aggregates
      per source IP to capture low-and-slow probing patterns. \\
  \addlinespace
  Feature Engineering
    & Time-window aggregation: distinct ports/IP/10s, SYN-ratio, IAT statistics, honeypot
      interaction count, and MTD port delta over rolling windows. \\
  \addlinespace
  Imbalance Handling
    & SMOTE oversampling on the \textsc{Port\_Scan} minority class in the training set. \\
  \addlinespace
  Evaluation Metrics
    & Accuracy, Precision, Recall, F1-Score, AUC-ROC, Detection Rate per scan type,
      Honeypot True Positive Rate. \\
  \addlinespace
  Validation Strategy
    & Stratified K-Fold ($k=10$) + temporal holdout (last 20\% by time order). \\
  \bottomrule
\end{tabular}
\label{tab:ml-strategy}
\end{table}

% ------------------------------------------------------------
\chapter{MTD \& Honeypot Component Specifications}

\section{Moving Target Defence (MTD)}

The MTD component continuously mutates key network configuration parameters to invalidate
any reconnaissance data an adversary may have collected. The core principle is to impose a
high and continuous cost on the attacker: a network map obtained through scanning becomes
stale within minutes, forcing repeated re-scanning that is itself detectable.

\begin{table}[h!]
\centering
\caption{MTD Component Specifications}
\begin{tabular}{@{}L{3.5cm} L{11cm}@{}}
  \toprule
  \textbf{Parameter} & \textbf{Details} \\
  \midrule
  Mutation Scope
    & Dynamic rotation of service port assignments and, optionally, virtual IP addresses
      of protected hosts within the network. \\
  \addlinespace
  Rotation Trigger
    & Time-based (scheduled rotation interval) and event-based (triggered automatically
      upon confirmed scanning detection by the AI engine). \\
  \addlinespace
  AI Integration
    & The MTD state table is fed as a live input to the AI model, ensuring that traffic
      directed at recently vacated ports is correctly identified as post-scan probing rather
      than legitimate communication. \\
  \addlinespace
  Coordination
    & All legitimate nodes on the network are synchronised with the current MTD state
      via a secure distribution mechanism, ensuring zero disruption to authorised traffic
      during rotation cycles. \\
  \bottomrule
\end{tabular}
\label{tab:mtd-specs}
\end{table}

\section{Honeypot Integration}

Honeypot nodes are strategically deployed within the network topology as high-fidelity decoy
systems. They present a realistic appearance to any scanning adversary while containing no
operational data or functionality. All traffic directed to a honeypot is inherently unauthorised
and constitutes an unambiguous threat signal.

\begin{table}[h!]
\centering
\caption{Honeypot Component Specifications}
\begin{tabular}{@{}L{3.5cm} L{11cm}@{}}
  \toprule
  \textbf{Parameter} & \textbf{Details} \\
  \midrule
  Honeypot Type
    & Low-to-medium interaction honeypots emulating common services (SSH, HTTP, FTP)
      to attract and engage scanning and post-scan probing activity. \\
  \addlinespace
  Placement Strategy
    & Distributed across the network topology to maximise the probability of contact
      during any broad or targeted scanning campaign. \\
  \addlinespace
  AI Contribution
    & Honeypot interaction events --- connection attempts, payloads, source IPs --- are
      ingested as real-time features by the AI model, significantly boosting detection
      confidence for confirmed threat actors. \\
  \addlinespace
  Profiling Output
    & Each honeypot interaction is logged and associated with the corresponding attacker
      profile, enriching the threat intelligence record with behavioural indicators and
      session-level metadata. \\
  \bottomrule
\end{tabular}
\label{tab:honeypot-specs}
\end{table}

% ------------------------------------------------------------
\chapter{Non-Functional Requirements}

\begin{table}[h!]
\centering
\caption{Non-Functional Requirements \& Performance Targets}
\begin{tabular}{@{}L{4.5cm} L{6cm} C{3.5cm}@{}}
  \toprule
  \textbf{Requirement} & \textbf{Criterion} & \textbf{Target Value} \\
  \midrule
  Fast Scan Detection    & Detection time for aggressive Nmap scan                   & $< 8$ seconds \\
  Slow Scan Detection    & Detection time for stealthy / slow scan                   & $< 90$\,s \\
  Model F1-Score         & On CICIDS2017 PortScan test subset                        & $\geq 88\%$ \\
  Model Accuracy         & Overall classification on test set                        & $\geq 90\%$ \\
  False Positive Rate    & Normal connections flagged as scans                       & $< 6\%$ \\
  Dashboard Refresh      & Scanner activity \& honeypot heat map update frequency    & $\leq 5$ seconds \\
  Blocking Response      & From detection to automated block rule enforcement        & $< 15$ seconds \\
  MTD Rotation Latency   & From detection trigger to completed port rotation         & $< 30$ seconds \\
  Honeypot Detection Rate & Proportion of scanning IPs interacting with a honeypot  & $\geq 60\%$ \\
  \bottomrule
\end{tabular}
\label{tab:nfr}
\end{table}

% ------------------------------------------------------------
\chapter{Project Phases \& Deliverables}

\begin{table}[h!]
\centering
\caption{Project Phases \& Deliverables}
\begin{tabular}{@{}L{3.5cm} L{6cm} L{5cm}@{}}
  \toprule
  \textbf{Phase} & \textbf{Activities} & \textbf{Key Deliverables} \\
  \midrule
  1 --- Conception
    & Network topology design, MTD architecture, feature set definition, scan type taxonomy.
    & Network diagram, MTD design spec, feature dictionary. \\
  \addlinespace
  2 --- Implementation
    & Dataset preprocessing, RF + XGBoost + Isolation Forest training, real-time capture
      pipeline, MTD rotation engine, honeypot deployment, dashboard.
    & Trained models (.pkl), IDS pipeline, MTD engine, honeypot nodes, dashboard UI. \\
  \addlinespace
  3 --- Testing
    & Simulate Nmap SYN, Connect, UDP, XMAS, slow scans; validate MTD disruption;
      test honeypot engagement.
    & Per-scan detection table, ROC curves, MTD effectiveness report, honeypot
      engagement log. \\
  \addlinespace
  4 --- Validation
    & K-Fold cross-validation, temporal holdout, end-to-end system validation including
      MTD and honeypot integration.
    & Final report, demo recording, multi-model comparison table. \\
  \bottomrule
\end{tabular}
\label{tab:phases}
\end{table}

% ------------------------------------------------------------
\chapter{Technical Stack Summary}

\begin{table}[h!]
\centering
\caption{Technical Stack Summary}
\begin{tabular}{@{}L{4cm} L{11cm}@{}}
  \toprule
  \textbf{Component} & \textbf{Technologies} \\
  \midrule
  OS / Virtualisation
    & Ubuntu Server (IDS + victim nodes), Kali Linux (attacker), GNS3 / VirtualBox \\
  Programming Language    & Python 3.10+ \\
  Traffic Capture
    & Scapy with BPF filters (SYN-only), PyShark, sliding window aggregation \\
  ML Libraries
    & Scikit-learn, XGBoost, imbalanced-learn, Joblib, Pandas, NumPy \\
  Web Dashboard
    & Flask + Socket.IO + D3.js (port heat map, honeypot event feed) / Grafana + InfluxDB \\
  Blocking Engine
    & Linux iptables DROP rules via Python subprocess, with TTL-based auto-unblock scheduler \\
  MTD Engine
    & Custom Python service managing dynamic port/IP rotation table, state distribution,
      and AI engine synchronisation \\
  Honeypot Framework
    & Cowrie (SSH/Telnet), Dionaea (multi-service), or custom lightweight Python socket
      listeners \\
  Attack Simulation
    & Nmap (SYN, ACK, UDP, XMAS, slow scans), Masscan, custom Scapy scripts \\
  Version Control         & Git + GitHub --- feature branches per module component \\
  \bottomrule
\end{tabular}
\label{tab:stack}
\end{table}

% ------------------------------------------------------------
\chapter{Expected Outcomes \& Success Criteria}

\begin{itemize}[nosep, leftmargin=2em, label=$\checkmark$]
  \item The trained model achieves an F1-Score $\geq 88\%$ on the CICIDS2017 PortScan
        test subset.
  \item An aggressive Nmap SYN stealth scan is detected and an alert is triggered within
        8 seconds.
  \item A slow-rate scan (1 probe/second) is correctly identified within the 90-second
        analysis window.
  \item The MTD engine completes a port rotation cycle within 30 seconds of a confirmed
        detection trigger.
  \item Following MTD rotation, continued probing of vacated ports by the same source IP
        is correctly classified as post-scan activity with no false negative.
  \item At least 60\% of scanning source IPs interact with at least one honeypot node
        during a full scan campaign, generating enriched attacker profile data.
  \item The dashboard reflects scanning events, honeypot interactions, and MTD rotation
        status with a maximum update delay of 5 seconds.
  \item Automated blocking is enforced within 15 seconds of a confirmed detection.
  \item The attacker profile correctly identifies scan type (SYN / UDP / ACK / XMAS) and
        includes honeypot contact history where applicable.
  \item Legitimate traffic from authorised nodes generates a false positive rate below 6\%
        across all test scenarios, including during active MTD rotation cycles.
\end{itemize}

% ============================================================
%  PART III — ARCHITECTURE & CONCEPTION
% ============================================================
\part{Architecture \& Conception --- Aegis Entropy}

\textit{Architecture Upgrade Note: The original Needs Expression outlined a standard defensive
posture. The final Aegis Entropy implementation elevates this to a proactive, highly aggressive
defence system. This part documents every design upgrade, its technical rationale, and the
three core pillars of the architecture.}

% ------------------------------------------------------------
\chapter{Architecture Upgrades: Blueprint vs.\ Reality}

The table below maps every original design choice from the Needs Expression document to its
upgraded counterpart in the final Aegis Entropy system, along with the technical rationale for
each upgrade.

\begin{table}[h!]
\centering
\caption{Architecture Upgrades: Original Specification vs.\ Aegis Entropy Implementation}
\begin{tabular}{@{}L{3cm} L{5cm} L{6.5cm}@{}}
  \toprule
  \textbf{Component} & \textbf{Original Specification} & \textbf{Aegis Entropy Upgrade} \\
  \midrule
  Response Engine
    & Linux iptables DROP rules via Python subprocess. Passive approach: the attacker's
      packets are silently discarded.
    & TCP Tarpit. Forged SYN-ACK with Window Size = 0 is returned, freezing tools and
      exhausting attacker resources. \\
  \addlinespace
  MTD Mechanism
    & Dynamic rotation of service port assignments and virtual IP addresses.
    & Active Fingerprint Mutation. Traffic intercepted inline; TCP/IP headers rewritten
      to falsify OS identity. \\
  \addlinespace
  Deception Layer
    & Full standalone honeypot VMs: Cowrie, Dionaea. High overhead.
    & Inline Shadow Nodes built with Python + Scapy. Fake responses generated directly
      from the network edge. \\
  \bottomrule
\end{tabular}
\label{tab:upgrades}
\end{table}

\textbf{Design Philosophy:} The shift from passive (DROP) to active (Tarpit + Mutation +
Shadow Nodes) defence reflects a deliberate strategic choice: instead of simply blocking the
attacker, Aegis Entropy wastes the attacker's time, corrupts their intelligence, and deceives
them --- all while the AI model continues to profile and classify the threat.

% ------------------------------------------------------------
\chapter{Pillar 1 --- Core Infrastructure: Asynchronous Interception via NFQUEUE}

The foundation of the Aegis Entropy active defence pipeline is the Linux kernel's
\textbf{Netfilter Queue (NFQUEUE)} mechanism. This provides a user-space hook into the
kernel's packet processing path, enabling Python scripts to intercept, inspect, and modify
--- or drop --- individual network packets in real time.

\begin{table}[h!]
\centering
\caption{NFQUEUE Interception Mechanism}
\begin{tabular}{@{}L{3cm} L{11.5cm}@{}}
  \toprule
  \textbf{Aspect} & \textbf{Description} \\
  \midrule
  Mechanism
    & When a packet reaches the NIC, the Linux kernel pauses its journey and hands it
      directly to the Python script via NFQUEUE. \\
  \addlinespace
  Verdict Power
    & The packet cannot proceed until the script issues one of three verdicts: ACCEPT,
      DROP, or MODIFY. \\
  \addlinespace
  Integration Point
    & All subsequent active defence mechanisms (Tarpit, Fingerprint Mutation, Shadow
      Node responses) are implemented as verdict decisions within this interception layer. \\
  \addlinespace
  Performance
    & Asynchronous processing ensures that the main network stack is not blocked; only
      flagged packets are diverted to the Python handler. \\
  \bottomrule
\end{tabular}
\label{tab:nfqueue}
\end{table}

% ------------------------------------------------------------
\chapter{Pillar 2 --- Sabotage Engine: TCP Tarpitting}

When the AI model flags an active Nmap scan, the Sabotage Engine engages: Scapy
intercepts the scanning probe and crafts a forged SYN-ACK response with Window Size = 0.
This is the core mechanism of \textbf{TCP Tarpitting}.

\section{Technical Mechanism}

A standard TCP connection requires both parties to negotiate a window size --- the amount of
data each end is willing to buffer. By returning Window Size = 0, the server signals that its
buffer is full and the client must wait before sending data. The attacker's machine and
scanning tools interpret this as: ``The server is alive but its buffer is full --- I must wait.''
Nmap will literally hang, frozen in an ESTABLISHED state, consuming one of its finite
connection threads.

\section{Operational Impact}

\begin{itemize}[nosep, leftmargin=1.5em]
  \item \textbf{Tool Freezing:} Nmap and similar scanners allocate a connection thread per
        probe. With Window Size = 0, the thread is blocked indefinitely, exhausting the
        scanner's parallelism.
  \item \textbf{Time Wasting:} A scan that would normally complete in seconds is stretched
        to minutes or hours.
  \item \textbf{Continued Profiling:} While the attacker is frozen, the AI model continues
        to accumulate behavioural data, refining the attacker profile and confidence score.
  \item \textbf{Deception:} The attacker receives a valid TCP response, confirming the port
        is open --- but receives no actual service data.
\end{itemize}

% ------------------------------------------------------------
\chapter{Pillar 3 --- Deception Engine: Active Fingerprint Mutation}

The \texttt{network\_mutator.py} script implements the \textbf{Active Fingerprint Mutation}
mechanism. Rather than rotating ports (the original Aegis-Morph approach), this upgrade
intercepts outgoing packets and dynamically rewrites TCP/IP header fields before they reach
the attacker, creating conflicting and inconsistent OS signatures from a single virtual machine.

\section{Mutation Scope \& Parameters}

\begin{table}[h!]
\centering
\caption{Active Fingerprint Mutation Parameters}
\begin{tabular}{@{}L{4cm} L{5.5cm} L{5cm}@{}}
  \toprule
  \textbf{Parameter} & \textbf{Description} & \textbf{Effect on Attacker} \\
  \midrule
  TTL Value
    & Dynamic rewriting of IP Time-to-Live field
    & Each probe response appears to originate from a different OS / hop count \\
  \addlinespace
  TCP Options --- MSS
    & Maximum Segment Size field in TCP options
    & Fingerprinting tools receive contradictory OS stack signatures \\
  \addlinespace
  TCP Options --- WScale
    & Window Scale factor in TCP handshake
    & Inconsistent with any single known OS profile \\
  \addlinespace
  TCP Options --- SACK
    & Selective Acknowledgement capability flag
    & Creates ambiguity in OS detection algorithms \\
  \addlinespace
  TCP Options --- Timestamps
    & Order and presence of TCP timestamp options
    & Disrupts timing-based OS fingerprinting (e.g., Nmap \texttt{-O}) \\
  \bottomrule
\end{tabular}
\label{tab:mutation-params}
\end{table}

\section{Rotation Triggers}

Fingerprint mutation is not static --- it operates under two complementary trigger modes:

\begin{itemize}[nosep, leftmargin=1.5em]
  \item \textbf{Time-Based Trigger:} Mutation parameters rotate on a scheduled interval,
        ensuring that even passive observers receive inconsistent fingerprint data over time.
  \item \textbf{Event-Based Trigger:} Upon a confirmed scan detection by the AI model, an
        immediate mutation cycle is initiated, invalidating any fingerprint data the attacker
        may have already collected.
\end{itemize}

\section{AI Integration}

The mutation state table is fed as a live input feature to the AI model. This ensures that
traffic directed at fingerprinted ports is correctly classified in context --- the model knows
which OS profile is currently being projected and can interpret probe responses accordingly.
This bidirectional integration between the deception layer and the detection engine is a key
architectural innovation of Aegis Entropy.

\bigskip
\noindent\textbf{Inline Shadow Nodes:} In addition to the three pillars above, Aegis Entropy
replaces full honeypot VMs with lightweight \textbf{Inline Shadow Nodes} --- Python socket
listeners combined with Scapy packet crafting --- that generate fake service responses (SSH,
HTTP, FTP banners) directly from the network edge. This eliminates the overhead of full OS
simulation while maintaining high-fidelity deception at the packet level.

% ============================================================
%  PART IV — ML DATA DICTIONARY
% ============================================================
\part{ML Data Dictionary}

% ------------------------------------------------------------
\chapter{Network Features Data Dictionary}

The following table details all network features extracted and used by the Machine Learning
model (Random Forest / XGBoost / Isolation Forest) for detection and active response.
Features marked \textit{Updated} or \textit{New} reflect the structural upgrades introduced by
Aegis Entropy.

\begin{longtable}{@{}L{4cm} C{2cm} L{8.5cm}@{}}
  \caption{Complete ML Feature Data Dictionary (Aegis Entropy Final Version)}
  \label{tab:data-dict} \\
  \toprule
  \textbf{Feature Name} & \textbf{Type} & \textbf{Description \& Relevance} \\
  \midrule
  \endfirsthead
  \toprule
  \textbf{Feature Name} & \textbf{Type} & \textbf{Description \& Relevance} \\
  \midrule
  \endhead
  \bottomrule
  \endfoot

  Distinct Dst Ports/IP
    & Numeric
    & Number of unique destination ports contacted by the source.
      \textit{Relevance: Critical --- core scanning signature.} \\
  \addlinespace
  SYN Flag Count
    & Numeric
    & SYN packets sent without completing the TCP handshake.
      \textit{Relevance: Critical --- stealth SYN scan indicator.} \\
  \addlinespace
  RST Flag Count
    & Numeric
    & RST responses received from closed ports.
      \textit{Relevance: Very High --- port probe rejection indicator.} \\
  \addlinespace
  Flow Duration
    & Numeric
    & Duration of each individual probe connection.
      \textit{Relevance: High --- usually very short in fast scans.} \\
  \addlinespace
  Total Fwd Packets
    & Numeric
    & Total packets per flow originating from the scanner.
      \textit{Relevance: High --- usually 1 packet per probe.} \\
  \addlinespace
  ACK Flag Count
    & Numeric
    & ACK packets sent without a prior SYN.
      \textit{Relevance: Medium --- ACK scan detection.} \\
  \addlinespace
  IAT Mean
    & Numeric
    & Mean inter-arrival time between successive probes.
      \textit{Relevance: High --- slow scans exhibit large IAT values.} \\
  \addlinespace
  Unique Dst IPs/Src
    & Numeric
    & Number of distinct destination hosts contacted by the source.
      \textit{Relevance: Medium --- indicates host discovery phase.} \\
  \addlinespace
  TTL Value
    & Numeric
    & Time-to-Live field in the IP header.
      \textit{Relevance: Low --- used in OS fingerprinting attempts.} \\
  \addlinespace
  Bwd Packet Length
    & Numeric
    & Response packet size (backwards direction).
      \textit{Relevance: Medium --- differentiates RST vs SYN-ACK ratio.} \\
  \addlinespace
  Shadow Node Interaction \textit{(Updated)}
    & Binary (Numeric)
    & Indicates whether the source IP contacted a decoy Inline Shadow Node.
      Replaces the original Honeypot Interaction Flag.
      \textit{Relevance: CRITICAL --- high-confidence threat signal.} \\
  \addlinespace
  MTD Port Delta
    & Numeric
    & Difference between the probed port number and the currently assigned port in
      the MTD rotation table.
      \textit{Relevance: High --- reveals the use of stale reconnaissance data.} \\
  \addlinespace
  TCP Window Size (outbound) \textit{(New)}
    & Numeric
    & Window size value in the SYN-ACK packet sent back to the attacker.
      \textit{Relevance: New --- metric introduced to evaluate TCP Tarpit effectiveness
      (Window Size = 0 signals active tarpitting).} \\
\end{longtable}

\section{Data Dictionary Legend}

\begin{table}[h!]
\centering
\caption{Data Dictionary Legend}
\begin{tabular}{@{}L{4cm} L{10.5cm}@{}}
  \toprule
  \textbf{Symbol / Highlight} & \textbf{Meaning} \\
  \midrule
  \textit{(Updated)}
    & Feature definition or implementation updated in the Aegis Entropy upgrade from
      the original Needs Expression. \\
  \addlinespace
  \textit{(New)}
    & Feature introduced for the first time in the Aegis Entropy implementation, not
      present in the original specification. \\
  \addlinespace
  CRITICAL relevance
    & Feature carries the highest predictive weight in the ensemble model; its presence
      or value is near-deterministic for threat classification. \\
  \bottomrule
\end{tabular}
\label{tab:legend}
\end{table}

\section{Model Usage Summary}

The 13 features above are consumed by the ensemble pipeline as follows:

\begin{itemize}[nosep, leftmargin=1.5em]
  \item \textbf{Random Forest and XGBoost} operate on the full feature vector for
        per-flow binary classification (\textsc{Benign} vs \textsc{Port\_Scan}).
  \item \textbf{Isolation Forest} operates on time-windowed aggregates (IAT Mean,
        Distinct Dst Ports/IP, SYN Flag Count, MTD Port Delta) for slow-scan anomaly
        detection.
  \item \textbf{Shadow Node Interaction} and \textbf{MTD Port Delta} serve as
        high-confidence override signals that can elevate a borderline
        \textsc{Benign} classification to \textsc{Port\_Scan} regardless of other feature
        values.
\end{itemize}

\vspace{2em}
\noindent\rule{\linewidth}{0.4pt}
\begin{center}
  \small Aegis Entropy --- Academic Project Report ---
  Module: Artificial Intelligence --- ENSA K\'{e}nitra --- 2025--2026
\end{center}

\end{document}